% Part: normal-modal-logic
% Chapter: tableaux
% Section: soundness

\documentclass[../../../include/open-logic-section]{subfiles}

\begin{document}

\olfileid[pt]{nml}{tab}{sou}

\olsection{Correção para \Log{K}}

\begin{editorial}
  Esta demonstração de correção reutiliza a demonstração de correção para
  a lógica proposicional clássica, O.s.\@, demonstra tudo desde o início.
  Isso é adequado quando se deseja uma demonstração autocontida. Para quem
  já viu a correção de tableaux comuns, haverá repetição. Pretende-se tornar
  possível alternar entre a versão autocontida e uma versão construída
  sobre o caso não modal.
\end{editorial}

\begin{explain}
  Para mostrar que os !!{tableau}s prefixados são corretos, precisamos
  mostrar que, se
  \[
  \sFmla{\True}{!B_1}[1], \dots, \sFmla{\True}{!B_n}[1], \sFmla{\False}{!A}[1]
  \]
  tem um !!{tableau} fechado, então $!B_1, \dots, !B_n \Entails !A$.
  É mais fácil demonstrar a contrapositiva: se, para algum $\mModel{M}$
  e algum mundo~$w$, $\mSat{M}{!B_i}[w]$ para todo $i=1$, \dots,~$n$,
  mas $\mSat/{M}{!A}[w]$, então nenhum !!{tableau} pode fechar. Tal
  contramodelo mostra que as hipóteses iniciais do !!{tableau} são
  satisfatíveis. A estratégia da demonstração é mostrar que, sempre que
  todas as !!{formula}s prefixadas em um ramo de !!a{tableau} são
  satisfatíveis, qualquer aplicação de uma regra produz ao menos um ramo
  estendido que também é satisfatível. Como ramos fechados são
  insatisfatíveis, todo !!{tableau} para um conjunto satisfatível de
  !!{formula}s prefixadas deve ter ao menos um ramo aberto.

  Para aplicar essa estratégia ao caso modal, precisamos estender nossa
  definição de ``satisfatível'' a modelos modais e prefixos. Feito isso,
  contudo, a demonstração é direta.
\end{explain}

\begin{defn}
  Seja $P$ algum conjunto de prefixos, O.s.\@, $P \subseteq (\PosInt)^*
  \setminus \{\emptyseq\}$, e seja $\mModel{M}$ um modelo. Uma função
  ~$f\colon P \to W$ é uma \emph{interpretação de}~$P$ em~$\mModel{M}$
  se, sempre que $\sigma$ e $\sigma.n$ pertencem ambos a~$P$, então
  $Rf(\sigma)f(\sigma.n)$.

  Em relação a uma interpretação dos prefixos~$P$, podemos definir:
  \begin{enumerate}
  \item $\mModel{M}$ satisfaz $\sFmla{\True}{!A}[\sigma]$ sse
    $\mSat{M}{!A}[f(\sigma)]$.
  \item $\mModel{M}$ satisfaz $\sFmla{\False}{!A}[\sigma]$ sse
    $\mSat/{M}{!A}[f(\sigma)]$.
  \end{enumerate}
\end{defn}

\begin{defn}
  Seja $\Gamma$ um conjunto de !!{formula}s prefixadas, e seja
  $P(\Gamma)$ o conjunto dos prefixos que nele ocorrem. Se $f$ é uma
  interpretação de~$P(\Gamma)$ em $\mModel{M}$, dizemos que $\mModel{M}$
  satisfaz $\Gamma$ em relação a~$f$, $\mSat{M}{\Gamma}[f]$, se
  $\mModel{M}$ satisfaz toda !!{formula} prefixada em~$\Gamma$ em relação
  a~$f$. $\Gamma$ é \emph{satisfatível} sse há um modelo~$\mModel{M}$ e
  uma interpretação~$f$ de $P(\Gamma)$ tais que
  $\mSat{M}{\Gamma}[f]$.
\end{defn}

\begin{prop}
  Se $\Gamma$ contém tanto $\sFmla{\True}{!A}[\sigma]$ quanto
  $\sFmla{\False}{!A}[\sigma]$, para alguma !!{formula}~$!A$ e algum
  prefixo~$\sigma$, então $\Gamma$ é insatisfatível.
\end{prop}

\begin{proof}
  Não pode haver um modelo~$\mModel{M}$ e uma interpretação~$f$ de
  $P(\Gamma)$ tais que valham tanto $\mSat{M}{!A}[f(\sigma)]$ quanto
  $\mSat/{M}{!A}[f(\sigma)]$.
\end{proof}

\begin{thm}[Correção]
  \ollabel{thm:tableau-soundness}
  Se $\Gamma$ tem um !!{tableau} fechado, $\Gamma$ é insatisfatível.
\end{thm}

\begin{proof}
Chamamos um ramo de !!a{tableau} de satisfatível sse o conjunto das
!!{signed formula}s nele é satisfatível, e chamamos !!a{tableau} de
satisfatível se contém ao menos um ramo satisfatível.

Mostramos o seguinte: estender um !!{tableau} satisfatível por uma das
regras de inferência sempre produz um !!{tableau} satisfatível. Isso
demonstrará o teorema: todo !!{tableau} fechado resulta da aplicação de
regras de inferência ao !!{tableau} que consiste apenas em hipóteses de
~$\Gamma$. Assim, se $\Gamma$ fosse satisfatível, todo !!{tableau} para
ele seria satisfatível. Um !!{tableau} fechado, porém, claramente não é
satisfatível, pois todos os seus ramos são fechados, e ramos fechados são
insatisfatíveis.

Suponha que tenhamos um !!{tableau} satisfatível, O.s.\@, !!a{tableau}
com ao menos um ramo satisfatível. Aplicar uma regra de inferência ou
acrescenta !!{signed formula}s a um ramo, ou divide um ramo em dois. Se
o !!{tableau} tem um ramo satisfatível que não é estendido pela aplicação
da regra em questão, ele permanece um ramo satisfatível no !!{tableau}
estendido; logo, o tableau estendido é satisfatível. Portanto, precisamos
considerar apenas o caso em que uma regra é aplicada a um ramo satisfatível.

Seja $\Gamma$ o conjunto das !!{signed formula}s nesse ramo, e seja
$\sFmla{S}{!A}[\sigma] \in \Gamma$ a !!{signed formula} à qual a regra
é aplicada. Se a regra não divide o ramo, precisamos mostrar que o ramo
estendido, O.s.\@, $\Gamma$ juntamente com as conclusões da regra, ainda
é satisfatível. Se a regra divide o ramo, precisamos mostrar que ao menos
um dos dois ramos resultantes é satisfatível.
\tagfalse{prvDiamond}
Primeiro, consideramos as inferências possíveis com apenas uma premissa.
\begin{enumerate}
\item O ramo é expandido pela aplicação de $\TRule{\True}{\lnot}$ a
  $\sFmla{\True}{\lnot !B}[\sigma] \in \Gamma$. O ramo estendido contém,
  então, as !!{signed formula}s $\Gamma \cup
  \{\sFmla{\False}{!B}[\sigma]\}$. Suponha $\mSat{M}{\Gamma}[f]$. Em
  particular, $\mSat{M}{\lnot !B}[f(\sigma)]$. Logo,
  $\mSat/{M}{!B}[f(\sigma)]$, O.s.\@, $\mModel{M}$ satisfaz
  $\sFmla{\False}{!B}[\sigma]$ em relação a~$f$.
\item O ramo é expandido pela aplicação de $\TRule{\False}{\lnot}$ a
  $\sFmla{\False}{\lnot !B}[\sigma] \in \Gamma$: exercício.
\item O ramo é expandido pela aplicação de $\TRule{\True}{\land}$ a
  $\sFmla{\True}{!B \land !C}[\sigma] \in \Gamma$, o que produz duas
  novas !!{signed formula}s no ramo: $\sFmla{\True}{!B}[\sigma]$ e
  $\sFmla{\True}{!C}[\sigma]$. Suponha $\mSat{M}{\Gamma}[f]$, em
  particular $\mSat{M}{!B \land !C}[f(\sigma)]$. Então
  $\mSat{M}{!B}[f(\sigma)]$ e $\mSat{M}{!C}[f(\sigma)]$. Isso significa
  que $\mModel{M}$ satisfaz tanto $\sFmla{\True}{!B}[\sigma]$ quanto
  $\sFmla{\True}{!C}[\sigma]$ em relação a~$f$.
\item O ramo é expandido pela aplicação de $\TRule{\False}{\lor}$ a
  $\sFmla{\False}{!B \lor !C} \in \Gamma$: exercício.
\item O ramo é expandido pela aplicação de $\TRule{\False}{\lif}$ a
  $\sFmla{\False}{!B \lif !C}[\sigma] \in \Gamma$: isso produz duas
  novas !!{signed formula}s no ramo, $\sFmla{\True}{!B}[\sigma]$ e
  $\sFmla{\False}{!C}[\sigma]$. Suponha $\mSat{M}{\Gamma}[f]$, em
  particular $\mSat/{M}{!B \lif !C}[f(\sigma)]$. Então
  $\mSat{M}{!B}[f(\sigma)]$ e $\mSat/{M}{!C}[f(\sigma)]$. Isso significa
  que $\mModel{M}, f$ satisfaz tanto $\sFmla{\True}{!B}[\sigma]$ quanto
  $\sFmla{\False}{!C}[\sigma]$.

\iftag{prvBox}{%
\item O ramo é expandido pela aplicação de $\TRule{\True}{\Box}$ a
  $\sFmla{\True}{\Box !B}[\sigma] \in \Gamma$:
  \iftag{probBox}{Exercício.}{Isso produz uma nova !!{signed
      formula}~$\sFmla{\True}{!B}[\sigma.n]$ no ramo, para algum
    $\sigma.n \in P(\Gamma)$ (pois $\sigma.n$ deve ser usado). Suponha
    $\mSat{M}{\Gamma}[f]$, em particular, $\mSat{M}{\Box
      !B}[f(\sigma)]$. Como $f$ é uma interpretação de prefixos e
    tanto $\sigma$ quanto $\sigma.n \in P(\Gamma)$, sabemos que
    $Rf(\sigma)f(\sigma.n)$. Logo, $\mSat{M}{!B}[f(\sigma.n)]$, O.s.\@,
    $\mModel{M}, f$ satisfaz $\sFmla{\True}{!B}[\sigma.n]$.}

\item O ramo é expandido pela aplicação de $\TRule{\False}{\Box}$ a
  $\sFmla{\False}{\Box !B}[\sigma] \in \Gamma$:
  \iftag{probBox}{Exercício.}{Isso produz uma nova !!{signed
      formula}~$\sFmla{\False}{!A}[\sigma.n]$, em que $\sigma.n$ é um
    prefixo novo no ramo, O.s.\@, $\sigma.n \notin P(\Gamma)$. Como
    $\Gamma$ é satisfatível, há uma $\Struct{M}$ e uma interpretação~$f$
    de $P(\Gamma)$ tais que $\Sat{M}{\Gamma}[f]$, em particular
    $\mSat/{M}{\Box !B}[f(\sigma)]$. Precisamos mostrar que
    $\Gamma \cup \{\sFmla{\False}{!B}[\sigma.n]\}$ é satisfatível. Para
    isso, definimos uma interpretação de $P(\Gamma) \cup \{\sigma.n\}$
    como segue:

  Como $\mSat/{M}{\Box !B}[f(\sigma)]$, há um $w \in W$ tal que
  $Rf(\sigma)w$ e $\mSat/{M}{!B}[w]$. Seja $f'$ como $f$, exceto que
  $f'(\sigma.n) = w$. Como $f'(\sigma) = f(\sigma)$ e $Rf(\sigma)w$,
  temos $Rf'(\sigma)f'(\sigma.n)$; logo, $f'$ é uma interpretação de
  $P(\Gamma) \cup \{\sigma.n\}$. Evidentemente,
  $\mSat/{M}{!B}[f'(\sigma.n)]$. Como $f(\sigma') = f'(\sigma')$ para
  todos os prefixos $\sigma' \in P(\Gamma)$, $\mSat{M}{\Gamma}[f']$.
  Portanto, $\mModel{M}, f'$ satisfaz $\Gamma \cup
  \{\sFmla{\False}{!B}[\sigma.n]\}$.}
  }{}

\iftag{prvDiamond}{%
\item O ramo é expandido pela aplicação de $\TRule{\True}{\Diamond}$ a
  $\sFmla{\True}{\Diamond !B}[\sigma] \in \Gamma$:
  \iftag{probDiamond}{Exercício.}{Isso produz uma nova !!{signed
      formula}~$\sFmla{\True}{!A}[\sigma.n]$, em que $\sigma.n$ é um
    prefixo novo no ramo, O.s.\@, $\sigma.n \notin P(\Gamma)$. Como
    $\Gamma$ é satisfatível, há uma $\Struct{M}$ e uma interpretação~$f$
    de $P(\Gamma)$ tais que $\Sat{M}{\Gamma}[f]$, em particular
    $\mSat{M}{\Diamond !B}[f(\sigma)]$. Precisamos mostrar que
    $\Gamma \cup \{\sFmla{\True}{!B}[\sigma.n]\}$ é satisfatível. Para
    isso, definimos uma interpretação de $P(\Gamma) \cup \{\sigma.n\}$
    como segue:

  Como $\mSat{M}{\Diamond !B}[f(\sigma)]$, há um $w \in W$ tal que
  $Rf(\sigma)w$ e $\mSat{M}{!B}[w]$. Seja $f'$ como $f$, exceto que
  $f'(\sigma.n) = w$. Como $f'(\sigma) = f(\sigma)$ e $Rf(\sigma)w$,
  temos $Rf'(\sigma)f'(\sigma.n)$; logo, $f'$ é uma interpretação de
  $P(\Gamma) \cup \{\sigma.n\}$. Evidentemente,
  $\mSat{M}{!B}[f'(\sigma.n)]$. Como $f(\sigma') = f'(\sigma')$ para
  todos os prefixos $\sigma' \in P(\Gamma)$, $\mSat{M}{\Gamma}[f']$.
  Portanto, $\mModel{M}, f'$ satisfaz $\Gamma \cup
  \{\sFmla{\True}{!B}[\sigma.n]\}$.}
\item O ramo é expandido pela aplicação de $\TRule{\False}{\Diamond}$ a
  $\sFmla{\False}{\Diamond !B}[\sigma] \in \Gamma$:
  \iftag{probDiamond}{Exercício.}{Isso produz uma nova !!{signed
      formula}~$\sFmla{\False}{!B}[\sigma.n]$ no ramo, para algum
    $\sigma.n \in P(\Gamma)$ (pois $\sigma.n$ deve ser usado). Suponha
    $\mSat{M}{\Gamma}[f]$, em particular, $\mSat/{M}{\Diamond
      !B}[f(\sigma)]$. Como $f$ é uma interpretação de prefixos e
    tanto $\sigma$ quanto $\sigma.n \in P(\Gamma)$, sabemos que
    $Rf(\sigma)f(\sigma.n)$. Logo, $\mSat/{M}{!B}[f(\sigma.n)]$, O.s.\@,
    $\mModel{M}, f$ satisfaz $\sFmla{\False}{!B}[\sigma.n]$.}
  }{}
\end{enumerate}
Consideremos agora as inferências possíveis com duas premissas.
\begin{enumerate}
\item O ramo é expandido pela aplicação de $\TRule{\False}{\land}$ a
  $\sFmla{\False}{!B \land !C}[\sigma] \in \Gamma$, o que produz dois
  ramos: um ramo esquerdo que continua por
  $\sFmla{\False}{!B}[\sigma]$ e um ramo direito que continua por
  $\sFmla{\False}{!C}[\sigma]$. Suponha $\mSat{M}{\Gamma}[f]$, em
  particular $\mSat/{M}{!B \land !C}[f(\sigma)]$. Então
  $\mSat/{M}{!B}[f(\sigma)]$ ou $\mSat/{M}{!C}[f(\sigma)]$. No primeiro
  caso, $\mModel{M}, f$ satisfaz $\sFmla{\False}{!B}[\sigma]$, O.s.\@,
  o ramo esquerdo é satisfatível. No segundo, $\mModel{M}, f$ satisfaz
  $\sFmla{\False}{!C}[\sigma]$, O.s.\@, o ramo direito é satisfatível.
\item O ramo é expandido pela aplicação de $\TRule{\True}{\lor}$ a
    $\sFmla{\True}{!B \lor !C}[\sigma] \in \Gamma$: exercício.
\item O ramo é expandido pela aplicação de $\TRule{\True}{\lif}$ a
    $\sFmla{\True}{!B \lif !C}[\sigma] \in \Gamma$: exercício.
\end{enumerate}
\end{proof}

\begin{prob}
Complete a demonstração de \olref[nml][tab][sou]{thm:tableau-soundness}.
\end{prob}

\begin{cor}
\ollabel{cor:entailment-soundness}
Se $\Gamma \Proves !A$, então $\Gamma \Entails !A$.
\end{cor}

\begin{proof}
  Se $\Gamma \Proves !A$, então, para alguns $!B_1$, \dots, $!B_n \in
  \Gamma$, $\Delta = \{\sFmla{\False}{!A}[1], \sFmla{\True}{!B_1}[1],
  \dots, \sFmla{\True}{!B_n}[1]\}$ tem um !!{tableau} fechado. Queremos
  mostrar que $\Gamma \Entails !A$. Suponha que não; então, para algum
  $\mModel{M}$ e algum $w$, $\mSat{M}{!B_i}[w]$ para $i=1$, \dots,~$n$,
  mas $\mSat/{M}{!A}[w]$. Seja $f(1) = w$; então $f$ é uma interpretação
  de~$P(\Delta)$ em~$\mModel{M}$, e $\mModel{M}$ satisfaz~$\Delta$ em
  relação a~$f$. Mas, por \olref{thm:tableau-soundness}, $\Delta$ é
  insatisfatível, pois tem um !!{tableau} fechado, uma contradição.
  Logo, afinal devemos ter $\Gamma \Proves !A$.
\end{proof}

\begin{cor}
\ollabel{cor:weak-soundness}
Se $\Proves !A$, então $!A$ é verdadeira em todos os modelos.
\end{cor}

\end{document}
